% Chapter 1

\chapter{Introdução} % Main chapter title
\label{chap:Chapter1} % For referencing the chapter elsewhere, use Chapter~\ref{Chapter1}

Este capítulo apresenta o contexto e a motivação do trabalho, a organização onde decorre e o problema abordado. Define depois os objetivos, as questões de investigação, a abordagem metodológica e a gestão de riscos. Termina com as contribuições esperadas, as considerações éticas e profissionais e a estrutura do documento.

%-------------------------------------------------------------------------------
\section{Contexto e Motivação}
\label{sec:chap1_context}

Os bancos de perguntas de escolha múltipla (\glspl{MCQ}) são um recurso habitual na preparação para exames na área da saúde. O Adhara QBank é uma plataforma de estudo da Adhara com perguntas de escolha múltipla em várias línguas. Estudantes de medicina, médicos e enfermeiros usam-na para preparar exames (Secção~\ref{sec:chap1_org}).

Durante o estudo, o utilizador pode ter dúvidas que a explicação da pergunta não resolve. Pode querer saber, por exemplo, porque é que uma opção está errada ou como um achado da vinheta se relaciona com a resposta correta. A plataforma não permite esclarecer essas dúvidas no momento. No entanto, a base de dados guarda mais informação sobre cada pergunta do que a que aparece no ecrã.


A literatura propõe assistentes conversacionais baseados em \glspl{LLM} como apoio ao estudo em educação e em educação médica \parencite{kasneci2023chatgpt,abdalrazaq2023llm}. Num ensaio controlado com estudantes de medicina, um destes assistentes trouxe benefícios para a aprendizagem \parencite{gan2024orthopedic}. Em contexto médico, porém, um assistente genérico traz riscos. Os \glspl{LLM} podem gerar respostas plausíveis mas incorretas ou infiéis ao contexto que recebem \parencite{hu2025advancing,jain2025aligning}, e essas respostas podem contradizer o conteúdo validado pela equipa científica. Fundamentar a resposta em informação dada ao modelo, por exemplo com \gls{RAG} \parencite{hu2025advancing}, reduz este risco \parencite{jang2025medtutor}, mas não o elimina \parencite{xie2026when,rahat2026noise}. Além disso, um utilizador pode contornar as instruções do assistente através da manipulação do \emph{prompt} \parencite{emekci2025securing,geng2026prompt}. Num exame, isso compromete a integridade da avaliação. A literatura aponta ainda a falta de orientações para o uso de \gls{IA} em educação médica \parencite{gordon2024scoping}.

%-------------------------------------------------------------------------------
\section{A Organização e o Adhara QBank}
\label{sec:chap1_org}

O trabalho decorre em colaboração com a Adhara (Adhara, Unipessoal, Lda.), uma empresa portuguesa de educação médica digital e de formação pós-graduada para profissionais de saúde. O Adhara QBank é a plataforma da empresa para a preparação de exames com perguntas de escolha múltipla. Existe em várias línguas e terá algumas centenas de utilizadores.

O utilizador resolve exames em três modos: Teste, Estudo e Simulado. A equipa científica da Adhara valida a informação de cada pergunta: vinheta clínica, opções de resposta, informação-chave, dicas, resposta correta, explicação, catalogação e informação complementar. Cada versão linguística do QBank tem o seu próprio conteúdo validado.

%-------------------------------------------------------------------------------
\section{Definição do Problema }
\label{sec:chap1_problem}

Ao resolver uma pergunta no Adhara QBank, o utilizador pode precisar de esclarecer aspetos que a explicação no ecrã não cobre. A base de dados guarda informação adicional sobre a pergunta, mas o utilizador não a consegue consultar com perguntas em linguagem natural. As fontes externas e os assistentes genéricos podem dar respostas que contradizem o conteúdo validado do item ou que acrescentam informação alheia ao contexto.

O problema desta dissertação é conceber um assistente conversacional baseado em \gls{LLM} que responda apenas com a informação que a base de dados guarda para cada pergunta. Essa informação inclui a vinheta, as opções, a catalogação, a resposta correta, a explicação, as dicas e a informação complementar. O assistente não deve recorrer a conhecimento externo e deve manter-se fiel a esse contexto. Também não deve responder a questões fora do âmbito da pergunta nem comprometer a integridade da avaliação.

O problema tem as seguintes características:
\begin{itemize}
    \item \textbf{Modos com regras diferentes.} O QBank tem modos de estudo e de avaliação. Uma resposta que revele a resposta correta antes de o utilizador responder anula o exercício num modo de avaliação.
    \item \textbf{Multilinguismo.} Cada versão linguística do QBank tem o seu próprio conteúdo validado. Uma resposta tem de ser coerente com o conteúdo da língua em que o utilizador estuda.
    \item \textbf{Contexto fechado por item.} O contexto relevante é conhecido à partida: são os dados da pergunta em curso. Não é preciso procurá-lo num corpus de documentos.
    \item \textbf{Restrições a harmonizar.} A solução tem de conciliar as capacidades conversacionais do modelo com a fidelidade ao contexto, a qualidade da resposta, a segurança, o custo, a latência e a integração na plataforma. Numa conversa, o utilizador espera uma resposta em poucos segundos, e cada pedido a um modelo comercial tem um custo por \emph{token}. Com algumas centenas de utilizadores e uma utilização desconhecida à partida, o custo e a latência de cada pedido são restrições de desenho e não apenas métricas de avaliação.
    \item \textbf{Várias partes interessadas.} Os utilizadores (estudantes, médicos e enfermeiros) e a Adhara, incluindo a equipa científica responsável pelo conteúdo.
    \item \textbf{Consequências.} Uma resposta fora do conteúdo validado pode induzir erros de aprendizagem e afetar a credibilidade da plataforma, a privacidade dos utilizadores e os custos de operação \parencite{abdalrazaq2023llm,hu2025advancing}.
\end{itemize}

%-------------------------------------------------------------------------------
\section{Objetivos}
\label{sec:chap1_objectives}

O objetivo geral é conceber, implementar e avaliar um assistente conversacional integrado no Adhara QBank. O assistente esclarece dúvidas sobre cada pergunta de escolha múltipla e usa apenas a informação da base de dados dessa pergunta.

Para atingir este objetivo, o trabalho define cinco objetivos específicos, cada um com uma evidência verificável:
\begin{description}
    \item[O1] Sistematizar a literatura sobre assistentes conversacionais baseados em \glspl{LLM} com respostas limitadas a um contexto definido, e sobre a avaliação dessas respostas. \emph{Evidência:} revisão com protocolo, diagrama \gls{PRISMA} e síntese por \gls{RQ} (Capítulo~\ref{chap:Chapter2}).
    \item[O2] Levantar com a Adhara os requisitos, as regras por modo e os campos da base de dados que formam o contexto de cada pergunta. \emph{Evidência:} requisitos e \emph{backlog} revistos pela Adhara.
    \item[O3] Implementar e integrar o assistente no fluxo de resolução de perguntas do QBank. \emph{Evidência:} versão integrada e demonstrada em perguntas reais.
    \item[O4] Aplicar as regras de cada modo e a recusa fora do âmbito, sem depender só das instruções ao modelo. \emph{Evidência:} testes de âmbito e adversariais aprovados.
    \item[O5] Avaliar a correção, a fidelidade, a recusa, o comportamento em várias línguas, a latência, o custo e a utilidade percebida. \emph{Evidência:} resultados da avaliação face aos critérios de sucesso definidos nos requisitos.
\end{description}

%-------------------------------------------------------------------------------
\section{Questões de Investigação}
\label{sec:chap1_rqs}

Duas questões de investigação (\glspl{RQ}) orientam o trabalho. A revisão sistemática do Capítulo~\ref{chap:Chapter2} responde-lhes, e as respostas orientam o desenho do assistente (objetivos O3 e O4) e a sua avaliação (objetivo O5):
\begin{description}
    \item[RQ1] Como utilizam os sistemas conversacionais baseados em \glspl{LLM} um contexto específico da aplicação para limitar as respostas, e que compromissos operacionais reportam essas abordagens?
    \item[RQ2] Como avaliam estes sistemas a correção factual, a fidelidade ao contexto fornecido e a adesão ao âmbito definido?
\end{description}

A RQ1 abrange a forma como o contexto chega ao modelo, os mecanismos de controlo do âmbito e de segurança e os compromissos de fiabilidade, latência, custo e integração. A RQ2 abrange os métodos, as métricas e os resultados da avaliação das respostas.

%-------------------------------------------------------------------------------
\section{Abordagem Metodológica}
\label{sec:chap1_method}

O trabalho segue a metodologia \gls{DSR}, que produz conhecimento através da conceção e avaliação de um artefacto que resolve um problema relevante \parencite{hevner2004design}. O trabalho segue as seis atividades propostas por \textcite{peffers2007design}:
\begin{enumerate}
    \item \textbf{Identificação do problema e motivação}: caracterização do problema com a Adhara (Secções~\ref{sec:chap1_context} e~\ref{sec:chap1_problem}).
    \item \textbf{Definição dos objetivos da solução}: objetivos e critérios de sucesso (Secção~\ref{sec:chap1_objectives}), apoiados na revisão da literatura (Capítulo~\ref{chap:Chapter2}).
    \item \textbf{Desenho e desenvolvimento}: arquitetura do assistente e integração no QBank, em iterações sobre um \emph{backlog} priorizado.
    \item \textbf{Demonstração}: utilização do assistente em perguntas reais do QBank.
    \item \textbf{Avaliação}: testes técnicos em várias línguas, avaliação por especialistas e questionário aos utilizadores.
    \item \textbf{Comunicação}: esta dissertação.
\end{enumerate}

Para construir o estado da arte, o autor conduziu uma revisão sistemática da literatura segundo as orientações de \textcite{kitchenham2007guidelines} e relatou-a segundo o \gls{PRISMA} 2020 \parencite{page2021prisma}. O Capítulo~\ref{chap:Chapter2} descreve o protocolo, a pesquisa, a seleção e a síntese.

%-------------------------------------------------------------------------------
\section{Gestão de Riscos}
\label{sec:chap1_risks}

O autor identificou os riscos do projeto no início do trabalho e classificou-os quanto à probabilidade (P) e ao impacto (I), numa escala de 1 (muito baixo) a 5 (muito alto). O nível de cada risco é o produto P\,$\times$\,I: baixo até 5, médio de 6 a 11 e alto a partir de 12. A Tabela~\ref{tab:risks} apresenta os riscos e as medidas de mitigação. Os riscos são revistos no fim de cada iteração do desenvolvimento.

\begin{table}[htbp]
\caption{Riscos do projeto, classificação e medidas de mitigação.}
\label{tab:risks}
\centering
\footnotesize
\begin{tabular}{p{0.5cm} p{4.2cm} c c c p{6.2cm}}
\toprule
\tabhead{Id} & \tabhead{Risco} & \tabhead{P} & \tabhead{I} & \tabhead{P$\times$I} & \tabhead{Mitigação} \\
\midrule
R1 & Respostas incorretas ou infiéis ao conteúdo validado da pergunta & 3 & 5 & 15 & Contexto fechado por pergunta; instruções de recusa; avaliação da fidelidade com perguntas de teste e por especialistas \\
R2 & Custo do modelo elevado, com uma utilização do assistente desconhecida à partida & 3 & 4 & 12 & Estimativa de custo por pedido; contexto compacto; cache de respostas; limites de utilização; monitorização do consumo \\
R3 & Tempo de resposta demasiado longo para uma conversa & 3 & 4 & 12 & Resposta em \emph{streaming}; contexto compacto; cache; medição da latência desde o início do desenvolvimento \\
R4 & Revelação da resposta correta ou desvio do âmbito por manipulação do \emph{prompt} & 3 & 4 & 12 & Regras de cada modo aplicadas fora do modelo (por exemplo, não enviar dados ainda não revelados); testes adversariais \\
R5 & Qualidade desigual entre as línguas do QBank & 3 & 3 & 9 & Conjunto de teste em todas as línguas; uso do conteúdo validado da língua do QBank \\
R6 & Pouca disponibilidade de especialistas para a avaliação & 3 & 3 & 9 & Planear a avaliação cedo; usar uma amostra; \gls{LLM} juiz como apoio, validado numa subamostra \\
R7 & Alteração de preço, de modelo ou indisponibilidade do fornecedor da \gls{API} & 2 & 4 & 8 & Camada de abstração do fornecedor; alternativa identificada; registo da versão do modelo usada \\
R8 & Atraso no cronograma & 3 & 3 & 9 & \emph{Backlog} priorizado; versão mínima funcional primeiro; acompanhamento com o orientador \\
R9 & Crescimento da base de dados com a cache e os registos & 2 & 2 & 4 & Política de retenção; exportação e limpeza periódica \\
\bottomrule
\end{tabular}
\end{table}

%-------------------------------------------------------------------------------
\section{Contribuições Esperadas}
\label{sec:chap1_contributions}

O trabalho pretende contribuir com:
\begin{itemize}
    \item uma revisão sistemática sobre assistentes conversacionais baseados em \glspl{LLM} com respostas limitadas a um contexto definido, em educação e saúde, com uma síntese dos mecanismos de contextualização, de controlo de âmbito e de avaliação (Capítulo~\ref{chap:Chapter2});
    \item um assistente conversacional integrado no fluxo de resolução de perguntas do Adhara QBank, que responde na língua do QBank com a informação validada da pergunta em curso e respeita as regras de cada modo;
    \item uma abordagem para entregar ao modelo o contexto fechado de cada pergunta e para recusar questões sem suporte nesse contexto;
    \item um protocolo e um conjunto de testes para avaliar a correção, a fidelidade ao contexto, a capacidade de recusa, o comportamento multilingue, a latência e o custo do assistente.
\end{itemize}

%-------------------------------------------------------------------------------
\section{Considerações Éticas e Profissionais}
\label{sec:chap1_ethics}

\paragraph{Código seguido.} Esta dissertação segue o Código Europeu de Conduta para a Integridade da Investigação \parencite{allea2023code}. O código assenta em quatro princípios: fiabilidade, honestidade, respeito e responsabilização. Considera inaceitáveis a fabricação, a falsificação e o plágio, e pede que os riscos da investigação e das suas aplicações sejam ponderados e que as incertezas dos resultados sejam reportadas. Os parágrafos seguintes mostram como estes princípios são aplicados.

\paragraph{Integridade da investigação.} O trabalho não fabrica, não falsifica nem omite dados ou resultados. Os resultados da avaliação são apresentados na íntegra, incluindo os negativos, com as incertezas e as limitações que restam. Todas as referências existem, foram lidas e foram verificadas na fonte original através do DOI. As figuras de outros autores são reproduzidas com a indicação da fonte e da licença. Para que a revisão possa ser repetida, as expressões de pesquisa, os registos exportados, as decisões de seleção com o respetivo motivo e os \emph{scripts} estão guardados e versionados (Apêndice~\ref{AppendixProtocol}).

\paragraph{Enviesamento.} A revisão foi feita por um único revisor, um risco de enviesamento discutido nas ameaças à validade (Secção~\ref{sec:chap2_threats}). Na avaliação do assistente, o conjunto de teste e os critérios de sucesso são definidos antes de se conhecerem os resultados, para evitar o enviesamento de confirmação. Os resultados são reportados por língua, para que nenhuma língua do QBank fique em desvantagem sem que isso seja visível.

\paragraph{Dissertação numa empresa.} O trabalho decorre ao abrigo de um protocolo de colaboração entre a Adhara e o ISEP. O conteúdo do QBank pertence à Adhara, é confidencial e o trabalho usa-o apenas neste projeto. A propriedade do código e dos resultados, e o que pode ser descrito na dissertação e em publicações, são acordados por escrito com a Adhara no início da dissertação. O autor desenvolve e avalia um produto para a empresa onde decorre o trabalho, o que constitui um potencial conflito de interesses. Para o mitigar, a avaliação segue um protocolo definido à partida, inclui especialistas e utilizadores, e os resultados desfavoráveis são reportados.

\paragraph{Participação humana e dados pessoais.} O questionário a utilizadores pede consentimento informado e explica o que é recolhido, porquê, durante quanto tempo e quem tem acesso. Recolhe apenas os dados necessários, guarda as respostas de forma anónima, em armazenamento institucional seguro, e nunca em contas pessoais nem em ferramentas de \gls{IA} externas. O tratamento segue o \gls{RGPD} \parencite{eu2016gdpr} e as orientações do ISEP. A necessidade de apreciação por uma comissão de ética para o questionário será confirmada com o orientador antes da recolha.

\paragraph{Impacto nos utilizadores.} Uma resposta errada pode induzir erros de aprendizagem num domínio em que esses erros têm consequências. As preocupações éticas traduzem-se em comportamento do sistema: o assistente usa apenas o contexto validado da pergunta, recusa o que está fora do âmbito, não revela a resposta correta antes do tempo e é testado contra utilizações abusivas, como a manipulação do \emph{prompt}. A correção das respostas é verificada por especialistas, porque o autor não tem formação médica. O assistente informa o utilizador de que interage com um sistema de \gls{IA}, como exige o Regulamento da Inteligência Artificial da União Europeia \parencite{eu2024aiact}.

\paragraph{Utilização de um modelo através de uma API externa.} O desenho compara um \gls{LLM} executado localmente com um \gls{LLM} acedido através de uma \gls{API} externa paga. Na segunda opção, o pedido do utilizador e o contexto da pergunta seguem para um serviço de terceiros. Para limitar os riscos:
\begin{itemize}
    \item o sistema envia apenas os dados necessários para responder: a mensagem do utilizador, o histórico da conversa sobre a pergunta e os dados dessa pergunta; não envia dados que identifiquem o utilizador;
    \item a escolha do fornecedor considera as suas condições de tratamento de dados, em particular a retenção dos pedidos, a não utilização para treino e a localização do processamento; a relação com o fornecedor segue o \gls{RGPD} \parencite{eu2016gdpr};
    \item o envio do conteúdo do QBank a um fornecedor externo exige a autorização da Adhara, detentora desse conteúdo.
\end{itemize}

\paragraph{Registo das perguntas para fins estatísticos.} O sistema poderá registar as perguntas feitas ao assistente e agrupá-las por tema, para fins estatísticos e de melhoria do conteúdo. Este registo terá uma finalidade definida, terá um prazo de conservação definido e será comunicado ao utilizador. Dados pseudonimizados continuam a ser dados pessoais. Por isso, ou o registo é anónimo, ou é tratado como dados pessoais segundo o \gls{RGPD}.
% TODO: confirmar com a Adhara a finalidade, a anonimização e o prazo de conservação.

\paragraph{Uso de ferramentas de IA generativa.} O autor é o único responsável pelo conteúdo desta dissertação; as ferramentas de \gls{IA} não são autoras. O autor usou ferramentas de \gls{IA} generativa nas seguintes tarefas:
\begin{itemize}
    \item \textbf{Revisão da literatura}: escrita dos \emph{scripts} de recolha, deduplicação e triagem; proposta de decisões de triagem por título e resumo; extração de dados dos textos integrais. O autor verificou cada decisão, que tem um motivo registado (Secção~\ref{sec:chap2_selection}), e os valores extraídos nos textos originais. As decisões finais de inclusão e a avaliação da qualidade são do autor.
    \item \textbf{Escrita}: apoio à estrutura e à revisão linguística do texto. O julgamento sobre o conteúdo e as conclusões é do autor.
    \item \textbf{Desenvolvimento de software}: escrita, revisão e depuração de código, verificado pelo autor com testes.
\end{itemize}
O conteúdo do QBank e o código da Adhara não são enviados a ferramentas de \gls{IA} externas sem a autorização da Adhara. O uso destas ferramentas é acordado com o orientador para cada tarefa. O próprio assistente usa um \gls{LLM} para gerar as respostas, como descrito acima.

%-------------------------------------------------------------------------------
\section{Estrutura do Documento}
\label{sec:chap1_structure}

Além deste capítulo, o documento tem a seguinte estrutura. O Capítulo~\ref{chap:Chapter2} apresenta o estado da arte: os fundamentos teóricos e uma revisão sistemática da literatura, com o protocolo, os resultados da pesquisa e da seleção, a síntese dos estudos e a lacuna de investigação. Os apêndices contêm o protocolo da revisão sistemática e a classificação dos estudos incluídos.
